We find that PHATE produces compelling visualizations with desirable features that simultaneously separate classes while preserving mesoscale and global connections. Distal points of the branches of the PHATE embeddings correspond to coarse-scale, archetypal classes, facilitating a straightforward manual identification. 

We find in all cases that PHATE produces more informative visual representations than PCA, UMAP, and t-SNE. We anticipate that this should inform more detailed and robust analyses of textual data, particularly in social science applications.

UMAP coupled with agglomerative clustering generally achieved the strongest performance across the clustering tasks, with consistently higher ARI scores across datasets. PHATE with agglomerative clustering also performed competitively, consistently outperforming PCA though slightly below UMAP on average in many cases. While PHATE visually preserves manifold structure better—particularly where UMAP compresses data into a line or single plane—the quantitative results support both UMAP and PHATE’s overall clustering effectiveness. It should be noted, however, that summary statistics across the clustering levels (i.e., the mean of the medians in Table \ref{tab:pivot_table}) capture only part of the story, and that, numerically, UMAP and PHATE achieve quite comparable results. In many cases, PHATE outperformed UMAP at coarser scales. With hyperparameter tuning, it is entirely possible that one could identify hyperparameter choices that would lead PHATE to outperform the reported results. 

Surprisingly, while satisfactorily recovering fine-scale clusterings in some cases, HDBSCAN consistently failed to produce adequate coarse-scale clusterings. The consistency of this finding suggests careful consideration is warranted when selecting a hierarchical clustering method.

This work focuses on PHATE as an alternative dimensionality reduction technique that emphasizes the preservation of both global and local data structure. Unlike UMAP, which prioritizes local neighborhood preservation and often fragments global trajectories, PHATE explicitly models the continuous diffusion of information across a data manifold, making it better suited for capturing smooth narrative arcs and multiscale semantic hierarchies. PHATE complements existing pipelines, like BERTopic, by offering a new embedding space in which hierarchical relationships and semantic trajectories are more naturally preserved. This approach can either serve as a drop-in replacement for dimensionality reduction in topic modeling workflows or augment existing methods by improving the interpretability and structure of visualizations. Compared to PCA, UMAP, and t-SNE, PHATE provides superior global semantic continuity, which is critical for uncovering latent hierarchies and narrative progressions in textual data. This could be helpful for measuring similarities and differences among graphs with nodes and edges decorated with textual features by utilizing the kernel-based approach suggested in \citet{berijanian2024soft}.

It could be interesting to apply PHATE in a sequential manner for hierarchical clustering in the following sense: Since the branches of the PHATE plots correspond to the most differentiated high-level classes, if one iteratively pruned these branches and reran PHATE on the remaining data, it might be possible to better-recover the nested classes. 
